\documentclass[12pt,a4paper]{article}

\usepackage[utf8]{inputenc}
\usepackage[T1]{fontenc}
\usepackage[spanish]{babel}
\usepackage{amsmath, amssymb}
\usepackage{graphicx}
\usepackage{booktabs}
\usepackage{hyperref}
\usepackage{geometry}
\usepackage{float}

\geometry{margin=2.5cm}

\title{\textbf{Optimizaci\'on de Rutas de Buses Escolares en Medell\'in}\\[6pt]
\large Comparaci\'on de Enfoques Heur\'isticos sobre Red Vial Real}
\author{Roy Sandoval \and Mariana \and Roberto C.\ Hinc\'apie \and Naomi \and Susana Toro}
\date{2026}

\begin{document}
\maketitle

\begin{abstract}
Este trabajo aborda el problema de ruteo de buses escolares con capacidad
en el \'area urbana de Medell\'in, Colombia.
Se propone una arquitectura de tres etapas basada en datos reales de la red vial:
generaci\'on del escenario con OpenStreetMap, clustering capacitado de estudiantes
y optimizaci\'on de rutas por cl\'uster.
Se implementan y comparan dos enfoques:
(1)~\textbf{K-medoids + Min-Cost Flow + TSP} (OR-Tools) y
(2)~\textbf{Set Cover via Programaci\'on Entera ILP} (CBC/PuLP).
\end{abstract}

\tableofcontents
\newpage

%-----------------------------------------------------------------------
\section{Introducci\'on}

El \textit{School Bus Routing Problem} (SBSRP) es una variante del Problema
de Ruteo de Veh\'iculos Capacitado (CVRP) donde el objetivo es planificar
rutas de buses para recoger estudiantes minimizando el costo total de
operaci\'on sujeto a la capacidad de cada bus.

En entornos urbanos, las distancias entre paradas dependen de la topolog\'ia
de la red vial.  Este trabajo usa datos reales de Medell\'in v\'ia OSMnx
para construir una matriz de distancias representativa del problema.

%-----------------------------------------------------------------------
\section{Descripci\'on del Problema}

\subsection{Instancia}

Dado un grafo dirigido ponderado $G=(V,E,w)$ que representa la red vial de Medell\'in:
\begin{itemize}
  \item $n=200$ estudiantes muestreados aleatoriamente sobre $V$ (semilla 42).
  \item Dep\'osito (colegio) $o\in V$ con coordenadas fijas.
  \item Capacidad m\'axima $C=20$ estudiantes por bus.
\end{itemize}

\subsection{Distancias por red vial}

Para cada par de nodos $i,j \in J\cup\{o\}$, la distancia real es el camino m\'as corto en $G$:
\begin{equation}
  d_{ij} = \min_{p\in\mathcal{P}(i,j)} \sum_{e\in p} w_e
\end{equation}
La matriz $D\in\mathbb{R}^{(n+1)\times(n+1)}$ resultante es generalmente asim\'etrica.

\subsection{Cota inferior de buses}

\begin{equation}
  k_{\min} = \left\lceil \frac{n}{C} \right\rceil = \left\lceil \frac{200}{20} \right\rceil = 10
\end{equation}

%-----------------------------------------------------------------------
\section{Generaci\'on del Escenario (\texttt{data\_generation.py})}

\begin{enumerate}
  \item Descarga de la red vial de Medell\'in con OSMnx;
        recorte a la regi\'on $6.182$--$6.312$\textdegree\,N,
        $75.531$--$75.626$\textdegree\,W.
  \item Extracci\'on de la componente fuertemente conexa (SCC) m\'as grande
        para garantizar alcanzabilidad entre todos los nodos.
  \item Muestreo de $n$ nodos de estudiantes y localizaci\'on del colegio.
  \item C\'alculo de la matriz de distancias mediante Dijkstra
        ($O(n|E|\log|V|)$).
\end{enumerate}

Para la fase de clustering, las distancias se simetrizan:
\begin{equation}
  D^{\mathrm{sym}}_{ij} = \frac{d_{ij}+d_{ji}}{2}
\end{equation}

%-----------------------------------------------------------------------
\section{Enfoque 1: K-medoids con Min-Cost Flow y TSP}

\subsection{Clustering capacitado}

El algoritmo itera entre dos fases hasta convergencia:

\paragraph{Asignaci\'on por flujo de costo m\'inimo.}
Dado el conjunto de medoides $M = \{m_1, \ldots, m_k\}$:
\begin{align}
  \min \quad & \sum_{i=1}^{n}\sum_{c=1}^{k} D^{\mathrm{sym}}_{i,m_c}\,x_{ic}
    \label{eq:km_obj}\\
  \text{s.t.}\quad
  & \sum_{c=1}^{k} x_{ic} = 1 \quad\forall i\in J \label{eq:km_assign}\\
  & \sum_{i=1}^{n} x_{ic} \leq C \quad\forall c \label{eq:km_cap}\\
  & x_{ic}\in\{0,1\} \label{eq:km_bin}
\end{align}
Este problema es equivalente a un flujo de costo m\'inimo y se resuelve en tiempo polinomial.

\paragraph{Actualizaci\'on de medoide.}
Para cada cl\'uster $S_c = \{i : x_{ic}=1\}$:
\begin{equation}
  m_c \leftarrow \arg\min_{h\in S_c}\sum_{i\in S_c} D^{\mathrm{sym}}_{ih}
\end{equation}

\subsection{Solver TSP por cl\'uster (\texttt{tsp\_solver.py})}

Para cada cl\'uster se define $V_c = S_c\cup\{o\}$ y se resuelve el TSP de
circuito cerrado con dep\'osito en $o$ usando OR-Tools
(\texttt{PATH\_CHEAPEST\_ARC} + \texttt{GUIDED\_LOCAL\_SEARCH}, l\'imite 5\,s):
\begin{equation}
  L_c = \min_{\sigma\in\mathcal{S}(S_c)}
        \left[ d_{o,\sigma_1}
             + \sum_{t=1}^{|S_c|-1} d_{\sigma_t,\sigma_{t+1}}
             + d_{\sigma_{|S_c|},o} \right]
\end{equation}
Distancia total de la flota: $L_{\mathrm{total}} = \sum_{c=1}^{k} L_c$.

%-----------------------------------------------------------------------
\section{Enfoque 2: Set Cover via Programaci\'on Entera}

\subsection{Pool de rutas candidatas}

Para cada estudiante $i\in J$ como centro, se construye una ruta greedy
con sus $C-1$ vecinos m\'as cercanos seg\'un $D^{\mathrm{sym}}$:
\begin{equation}
  r_i = \{i\}\cup
        \operatorname{top\text{-}(C-1)}\bigl\{j\in J\setminus\{i\}:
        D^{\mathrm{sym}}_{ij}\bigr\}
\end{equation}
El pool $\mathcal{R}=\{r_i\}_{i\in J}$ (sin duplicados) contiene
$m\leq n$ rutas con $|r_j|\leq C$.

\subsection{Formulaci\'on ILP}

Sea $x_j\in\{0,1\}$ la variable de selecci\'on de la ruta $j$:
\begin{align}
  \min \quad & \sum_{j=1}^{m} x_j \label{eq:sc_obj}\\
  \text{s.t.}\quad
  & \sum_{j:\,i\in r_j} x_j \geq 1 \quad\forall i\in J \label{eq:sc_cover}\\
  & x_j\in\{0,1\} \label{eq:sc_bin}
\end{align}
La restricci\'on~(\ref{eq:sc_cover}) garantiza que cada estudiante sea recogido
por al menos un bus.  El objetivo~(\ref{eq:sc_obj}) minimiza directamente el
n\'umero de buses.  Se resuelve con el solver CBC (PuLP) a optimalidad garantizada.

\subsection{Distancia de las rutas seleccionadas}

La distancia de cada ruta se aproxima con la heur\'istica del vecino m\'as
cercano (NN) partiendo del colegio $o$:
\begin{equation}
  \hat{L}_c = d_{o,\pi_1}
            + \sum_{t=1}^{|S_c|-1} d_{\pi_t,\pi_{t+1}}
            + d_{\pi_{|S_c|},o}
\end{equation}
donde $\pi$ es la permutaci\'on greedy NN.

%-----------------------------------------------------------------------
\section{Resultados}

\begin{table}[H]
  \centering
  \caption{Comparaci\'on de m\'etricas entre los dos enfoques
           (200 estudiantes, $C=20$, semilla~42).}
  \label{tab:comp}
  \begin{tabular}{lcc}
    \toprule
    \textbf{M\'etrica} & \textbf{K-medoids + TSP} & \textbf{Set Cover ILP}\\
    \midrule
    Buses utilizados         & 10    & 16     \\
    Distancia total (km)     & 261.19 & 455.08 \\
    Distancia media/bus (km) & 26.12  & 28.44  \\
    Tama\~no m\'in.\ (est.)  & 20    & 20     \\
    Tama\~no m\'ax.\ (est.)  & 20    & 20     \\
    Tama\~no promedio (est.) & 20.0  & 20.0   \\
    \bottomrule
  \end{tabular}\\[4pt]
  \small Escenario: 200 estudiantes, $C=20$, semilla~42. Valores obtenidos ejecutando la pipeline completa.
\end{table}

\begin{figure}[H]
  \centering
  \includegraphics[width=\linewidth]{comparacion_sizes.png}
  \caption{Distribuci\'on del n\'umero de estudiantes por bus en cada enfoque.
           La l\'inea roja indica la capacidad m\'axima $C=20$.}
  \label{fig:sizes}
\end{figure}

\begin{figure}[H]
  \centering
  \includegraphics[width=\linewidth]{comparacion_mapa.png}
  \caption{Asignaci\'on geogr\'afica de estudiantes a buses sobre el mapa de
           Medell\'in.  Cada color representa un bus distinto;
           la estrella roja es el colegio (dep\'osito).}
  \label{fig:maps}
\end{figure}

\subsection{Discusi\'on}

El enfoque \textbf{K-medoids + TSP} fija $k=k_{\min}$ antes de optimizar
y produce tours de alta calidad gracias a OR-Tools.
El enfoque \textbf{Set Cover ILP} aborda directamente el objetivo operacional
de minimizar la flota: la soluci\'on de cobertura es \'optima para el pool
generado, aunque los tours individuales son aproximados (heur\'istica NN).

%-----------------------------------------------------------------------
\section{Conclusiones}

\begin{itemize}
  \item Las distancias reales por red vial difieren significativamente de la
        distancia euclideana, lo que justifica el uso de OSMnx y Dijkstra.
  \item K-medoids + TSP produce tours de alta calidad con n\'umero de buses
        predeterminado por $k_{\min}$.
  \item Set Cover ILP minimiza directamente la flota con cobertura garantizada,
        a costa de tours aproximados.
  \item Ambos enfoques respetan estrictamente la capacidad $C=20$.
\end{itemize}

Como trabajo futuro se propone: combinar ambas etapas en un modelo unificado,
incorporar ventanas de tiempo (CVRPTW) y escalar a instancias con m\'ultiples colegios.

%-----------------------------------------------------------------------
\bibliographystyle{plain}
\begin{thebibliography}{9}

\bibitem{osmnx}
Boeing, G.\ (2017).
OSMnx: New methods for acquiring, constructing, analyzing, and visualizing
complex street networks.
\textit{Computers, Environment and Urban Systems}, 65, 126--139.

\bibitem{kmedoids}
Schubert, E.\ \& Rousseeuw, P.J.\ (2019).
Faster k-Medoids Clustering: Improving the PAM, CLARA, and CLARANS Algorithms.
\textit{SISAP 2019}, LNCS~11807.

\bibitem{ortools}
Google OR-Tools (2023).
\textit{Vehicle Routing}.
\url{https://developers.google.com/optimization/routing}

\bibitem{schittekat}
Schittekat, P., Kinable, J., Sorensen, K., Sevaux, M., Spieksma, F.,
\& Springael, J.\ (2013).
A metaheuristic for the school bus routing problem with bus stop selection.
\textit{European Journal of Operational Research}, 229(2), 518--528.

\bibitem{pulp}
Mitchell, S., OSullivan, M., \& Dunning, I.\ (2011).
PuLP: A Linear Programming Toolkit for Python.
The University of Auckland.

\end{thebibliography}

\end{document}
